\hspace{3mm}

\centerline{\bf \huge 7 класс}

\begin{flushright}
        \scriptsize{ $-$ Если сдаваться, не попробовав, результат не изменится.}
\end{flushright}


\problem{1}
Сугуру, Сукуна и Сатору решают олимпиаду ЭлМШ. Пока Сугуру решает $3$ задачи, Сукуна решает $4$, а пока Сукуна решает $3$ задачи, Сатору решает $5$ задач. Сколько задач решит Сугуру, если все втроём решили в общей сложности $9225$ задач?


\problem{2}
Директор Элистинского технического лицея взяла два трёхзначных числа, возвела каждое в квадрат и записала полученные числа друг за другом без пробела. Получилось некоторое десятизначное число. Могут ли не менее семи цифр в записи этого десятизначного числа равняться 3?

\problem{3}
В турнире по настольному теннису участвовали $2048$ человек. В каждом туре участники разбивались на пары и играли, после чего проигравшие выбывали из турнира (ничьих не бывает). После турнира, когда определился единоличный победитель, каждый из участников заявил, что в течение турнира обыграл не более одного правши. Причём все правши сказали правду, а все левши соврали. Сколько правшей могло участвовать в турнире?

\problem{4}
По кругу расставлено $90$ чисел, каждое из которых равно $1$, $2$ или $3$. Анджур Аркадьевич посчитал все $90$ сумм рядом стоящих чисел и обнаружил, что ровно $25$ из этих сумм равны $3$. Докажите, что какие-то два стоящих рядом числа равны.


\problem{5}
 Дана клетчатая доска 1000$\times$1000. Фигура \textit{гепард} из произвольной клетки $X$ бьёт все клетки квадрата 19$\times$19 с центральной клеткой $X$, за исключением клеток, находящихся с $X$ в одном столбце или одной строке. Какое наибольшее количество \textit{гепардов}, не бьющих друг друга, можно расставить на доске? (\textit{Богданов И.И.})